% Options for packages loaded elsewhere
% Options for packages loaded elsewhere
\PassOptionsToPackage{unicode}{hyperref}
\PassOptionsToPackage{hyphens}{url}
%
\documentclass[
  11pt,
  ignorenonframetext,
  aspectratio=169,
]{beamer}
\newif\ifbibliography
\usepackage{pgfpages}
\setbeamertemplate{caption}[numbered]
\setbeamertemplate{caption label separator}{: }
\setbeamercolor{caption name}{fg=normal text.fg}
\beamertemplatenavigationsymbolsempty
% remove section numbering
\setbeamertemplate{part page}{
  \centering
  \begin{beamercolorbox}[sep=16pt,center]{part title}
    \usebeamerfont{part title}\insertpart\par
  \end{beamercolorbox}
}
\setbeamertemplate{section page}{
  \centering
  \begin{beamercolorbox}[sep=12pt,center]{section title}
    \usebeamerfont{section title}\insertsection\par
  \end{beamercolorbox}
}
\setbeamertemplate{subsection page}{
  \centering
  \begin{beamercolorbox}[sep=8pt,center]{subsection title}
    \usebeamerfont{subsection title}\insertsubsection\par
  \end{beamercolorbox}
}
% Prevent slide breaks in the middle of a paragraph
\widowpenalties 1 10000
\raggedbottom
\AtBeginPart{
  \frame{\partpage}
}
\AtBeginSection{
  \ifbibliography
  \else
    \frame{\sectionpage}
  \fi
}
\AtBeginSubsection{
  \frame{\subsectionpage}
}
\usepackage{iftex}
\ifPDFTeX
  \usepackage[T1]{fontenc}
  \usepackage[utf8]{inputenc}
  \usepackage{textcomp} % provide euro and other symbols
\else % if luatex or xetex
  \usepackage{unicode-math} % this also loads fontspec
  \defaultfontfeatures{Scale=MatchLowercase}
  \defaultfontfeatures[\rmfamily]{Ligatures=TeX,Scale=1}
\fi
\usepackage{lmodern}

\ifPDFTeX\else
  % xetex/luatex font selection
\fi
% Use upquote if available, for straight quotes in verbatim environments
\IfFileExists{upquote.sty}{\usepackage{upquote}}{}
\IfFileExists{microtype.sty}{% use microtype if available
  \usepackage[]{microtype}
  \UseMicrotypeSet[protrusion]{basicmath} % disable protrusion for tt fonts
}{}
\makeatletter
\@ifundefined{KOMAClassName}{% if non-KOMA class
  \IfFileExists{parskip.sty}{%
    \usepackage{parskip}
  }{% else
    \setlength{\parindent}{0pt}
    \setlength{\parskip}{6pt plus 2pt minus 1pt}}
}{% if KOMA class
  \KOMAoptions{parskip=half}}
\makeatother


\usepackage{longtable,booktabs,array}
\newcounter{none} % for unnumbered tables
\usepackage{calc} % for calculating minipage widths
\usepackage{caption}
% Make caption package work with longtable
\makeatletter
\def\fnum@table{\tablename~\thetable}
\makeatother
\usepackage{graphicx}
\makeatletter
\newsavebox\pandoc@box
\newcommand*\pandocbounded[1]{% scales image to fit in text height/width
  \sbox\pandoc@box{#1}%
  \Gscale@div\@tempa{\textheight}{\dimexpr\ht\pandoc@box+\dp\pandoc@box\relax}%
  \Gscale@div\@tempb{\linewidth}{\wd\pandoc@box}%
  \ifdim\@tempb\p@<\@tempa\p@\let\@tempa\@tempb\fi% select the smaller of both
  \ifdim\@tempa\p@<\p@\scalebox{\@tempa}{\usebox\pandoc@box}%
  \else\usebox{\pandoc@box}%
  \fi%
}
% Set default figure placement to htbp
\def\fps@figure{htbp}
\makeatother





\setlength{\emergencystretch}{3em} % prevent overfull lines

\providecommand{\tightlist}{%
  \setlength{\itemsep}{0pt}\setlength{\parskip}{0pt}}



 


\setbeamertemplate{navigation symbols}{}
\setbeamertemplate{footline}[frame number]
\setbeamertemplate{itemize items}[circle]
\setbeamertemplate{itemize subitem}[triangle]
\setbeamercolor{frametitle}{fg=black}
\setbeamerfont{frametitle}{series=\bfseries}
\newcommand{\indep}{\perp\!\!\!\perp}
\usepackage{booktabs}
\usepackage{tikz}
\usetikzlibrary{arrows.meta,positioning,calc}
\tikzset{
  dnode/.style={circle,draw,thick,minimum size=9mm,inner sep=1pt,font=\small},
  unode/.style={circle,draw,thick,dashed,minimum size=9mm,inner sep=1pt,font=\small},
  bnode/.style={rectangle,draw,thick,fill=gray!25,minimum size=9mm,inner sep=3pt,font=\small},
  de/.style={-{Stealth[length=2.2mm]},thick},
  dde/.style={-{Stealth[length=2.2mm]},thick,dashed}
}
% shrink code input (fancyvrb) and code output (verbatim)
\usepackage{fancyvrb}
\fvset{fontsize=\scriptsize}
\makeatletter
\def\verbatim@font{\scriptsize\ttfamily}
\makeatother
\makeatletter
\@ifpackageloaded{caption}{}{\usepackage{caption}}
\AtBeginDocument{%
\ifdefined\contentsname
  \renewcommand*\contentsname{Table of contents}
\else
  \newcommand\contentsname{Table of contents}
\fi
\ifdefined\listfigurename
  \renewcommand*\listfigurename{List of Figures}
\else
  \newcommand\listfigurename{List of Figures}
\fi
\ifdefined\listtablename
  \renewcommand*\listtablename{List of Tables}
\else
  \newcommand\listtablename{List of Tables}
\fi
\ifdefined\figurename
  \renewcommand*\figurename{Figure}
\else
  \newcommand\figurename{Figure}
\fi
\ifdefined\tablename
  \renewcommand*\tablename{Table}
\else
  \newcommand\tablename{Table}
\fi
}
\@ifpackageloaded{float}{}{\usepackage{float}}
\floatstyle{ruled}
\@ifundefined{c@chapter}{\newfloat{codelisting}{h}{lop}}{\newfloat{codelisting}{h}{lop}[chapter]}
\floatname{codelisting}{Listing}
\newcommand*\listoflistings{\listof{codelisting}{List of Listings}}
\makeatother
\makeatletter
\makeatother
\makeatletter
\@ifpackageloaded{caption}{}{\usepackage{caption}}
\@ifpackageloaded{subcaption}{}{\usepackage{subcaption}}
\makeatother

\usepackage{bookmark}
\IfFileExists{xurl.sty}{\usepackage{xurl}}{} % add URL line breaks if available
\urlstyle{same}
\hypersetup{
  pdftitle={Unconfoundedness II: Propensity Scores and Regression},
  pdfauthor={Professor Ji-Woong Chung},
  hidelinks,
  pdfcreator={LaTeX via pandoc}}


\title{\texorpdfstring{Unconfoundedness II: Propensity Scores and
Regression}{Unconfoundedness II: Propensity Scores and Regression}}
\subtitle{\texorpdfstring{BUSS975: Causal Inference in Financial
Research}{BUSS975: Causal Inference in Financial Research}}
\author{Professor Ji-Woong Chung}
\date{}
\institute{Korea University}

\begin{document}
\frame{\titlepage}


\begin{frame}{Outline}
\protect\phantomsection\label{outline}
\begin{enumerate}
\tightlist
\item
  Propensity scores
\item
  Regression
\item
  The NSW demonstration
\end{enumerate}
\end{frame}

\section{1. Propensity Scores}\label{propensity-scores}

\begin{frame}{Where Lecture 3a left us}
\protect\phantomsection\label{where-lecture-3a-left-us}
\begin{itemize}
\tightlist
\item
  Unconfoundedness says: \textbf{compare treated and controls with the
  same \(X\).}
\item
  \textbf{Subclassification} and \textbf{matching} suffer from the curse
  of dimensionality.
\item
  Must we match on all of \(X\)? The assumption is stated conditional on
  \(X\), so it seems so.
\item
  \textbf{Rosenbaum and Rubin (1983): no.} One scalar computed from
  \(X\) does the whole job.
\end{itemize}
\end{frame}

\begin{frame}{Defining the propensity score}
\protect\phantomsection\label{defining-the-propensity-score}
\begin{itemize}
\tightlist
\item
  \textbf{Propensity score}: \(e(X) = \Pr(D = 1 \mid X)\)

  \begin{itemize}
  \tightlist
  \item
    ``the conditional probability of assignment to a particular
    treatment given a vector of observed covariates'' (Rosenbaum and
    Rubin 1983)
  \end{itemize}
\item
  Three things it is:

  \begin{itemize}
  \tightlist
  \item
    A \textbf{probability}, so it lives in \([0,1]\).
  \item
    A \textbf{conditional} probability --- read it as: \emph{of the
    units that look like this on \(X\), what fraction are treated?}
  \item
    A description of the \textbf{assignment mechanism}, not of the
    outcome.
  \end{itemize}
\item
  Two results follow: the \textbf{balancing property} (conditional on
  \(e(X)\), treated and control look alike on \(X\)) and the
  \textbf{propensity score theorem} (conditioning on \(e(X)\) does what
  conditioning on all of \(X\) does).
\end{itemize}
\end{frame}

\begin{frame}{The balancing property}
\protect\phantomsection\label{the-balancing-property}
\begin{center}
\fbox{\;\textbf{Claim.}\quad $D \indep X \mid e(X)$\;}
\end{center}

\scriptsize

\(D\) is binary, so its distribution is summarised by its mean; the
claim is equivalent to
\(E\big[D \mid X, e(X)\big] \;=\; E\big[D \mid e(X)\big]\). We show
\textbf{both sides equal \(e(X)\)}, and hence each other.

\textbf{Left side.} \(e(X)\) is a function of \(X\), so conditioning on
it adds nothing to \(X\): \[
E\big[D \mid X, e(X)\big] \;=\; E\big[D \mid X\big] \;=\; e(X)
\qquad \text{(definition of } e(X) \text{)} .
\] \textbf{Right side.} Apply the \textbf{law of iterated expectations},
\(E[A \mid C] = E_B\big[E[A \mid B, C] \mid C\big]\), with \(A = D\),
\(B = X\), \(C = e(X)\): \[
E\big[D \mid e(X)\big]
\;=\; E_{X}\Big[\;\underbrace{E\big[D \mid X, e(X)\big]}_{=\;e(X)\ \text{by the left side}} \;\Big|\; e(X)\Big]
\;=\; E_{X}\big[e(X) \mid e(X)\big] \;=\; e(X). \;\blacksquare
\]

\normalsize

\begin{itemize}
\tightlist
\item
  The mean version, \(E[X \mid D{=}1, e(X)] = E[X \mid D{=}0, e(X)]\),
  is what a balance table tests.
\end{itemize}
\end{frame}

\begin{frame}{The propensity score theorem}
\protect\phantomsection\label{the-propensity-score-theorem}
\begin{center}
\fbox{\;\textbf{Claim.}\quad If $(Y^1, Y^0) \indep D \mid X$, then $(Y^1, Y^0) \indep D \mid e(X)$\;}
\end{center}

\scriptsize

In mean form,
\(\underbrace{E\big[D \mid Y^1,Y^0, X\big] = E\big[D \mid X\big]}_{\text{assumed}} \quad \Longrightarrow \quad \underbrace{E\big[D \mid Y^1,Y^0, e(X)\big] = E\big[D \mid e(X)\big]}_{\text{to show}}\).

The right-hand side is \(e(X)\) by the previous slide, so only the
left-hand side is left to prove: \[
\begin{aligned}
E\big[D \mid Y^1,Y^0, e(X)\big]
&= E_{X}\Big[\, E\big[D \mid Y^1,Y^0, X, e(X)\big] \;\Big|\; Y^1,Y^0, e(X)\Big]
&& \text{law of iterated exp.} \\[1pt]
&= E_{X}\Big[\, E\big[D \mid Y^1,Y^0, X\big] \;\Big|\; Y^1,Y^0, e(X)\Big]
&& \text{$e(X)$ is a function of $X$} \\[1pt]
&= E_{X}\Big[\, \underbrace{E\big[D \mid X\big]}_{=\;e(X)} \;\Big|\; Y^1,Y^0, e(X)\Big] \;=\; e(X). \;\blacksquare
&& \textbf{unconfoundedness}
\end{aligned}
\]

\vfill

\normalsize

\begin{itemize}
\tightlist
\item
  The \(ATE\) is identified by comparing treated and control units at
  the same \textbf{score} rather than the same \(X\):
  \(ATE \;=\; E\Big[\, E\big[Y \mid D{=}1, e(X)\big] - E\big[Y \mid D{=}0, e(X)\big] \,\Big]\)
\end{itemize}
\end{frame}

\begin{frame}{Estimating the score}
\protect\phantomsection\label{estimating-the-score}
\begin{itemize}
\tightlist
\item
  You know the \textbf{true} score only if you ran the experiment.
  Otherwise estimate it: \(\Pr(D_i {=} 1 \mid X_i) = F(\beta X_i)\), by
  \textbf{logit or probit} --- not OLS, which extrapolates linearly and
  will hand you ``probabilities'' outside \([0,1]\).
\item
  Judge the specification by \textbf{balance achieved}, not by fit. A
  score with a high pseudo-\(R^2\) but poor balance is a bad score; the
  reverse is fine.
\item
  Fit, check balance, add interactions and polynomials where balance
  fails, refit. A legitimate specification search: it never looks at
  \(Y\), so the stage is \textbf{blind to the outcome}. Say so in the
  paper.
\item
  \textbf{The score creates no support that \(X\) lacks.} If a cell is
  all-treated or all-control, collapsing \(X\) to a scalar hides that;
  it does not fix it.
\end{itemize}
\end{frame}

\begin{frame}{Checking balance: not a \(t\)-test}
\protect\phantomsection\label{checking-balance-not-a-t-test}
The balancing property is checkable. Report it covariate by covariate,
before and after.

\[
\text{SMD} \;=\; \frac{\bar X_1 - \bar X_0}{\sqrt{\big(s_1^2 + s_0^2\big)/2}},
\qquad\qquad
\text{VR} \;=\; \frac{s_1^2}{s_0^2}
\]

\begin{itemize}
\tightlist
\item
  \textbf{Standardized mean difference.} \(|\text{SMD}| < 0.1\) is good.

  \begin{itemize}
  \tightlist
  \item
    Keep the denominator \textbf{fixed at its raw-sample value} when
    recomputing after adjustment, or the yardstick moves with the thing
    you are measuring.
  \end{itemize}
\item
  \textbf{Variance ratio.} Means can match while spreads do not.
  Rubin(2001): stay inside \([0.5,\,2]\), ideally \([0.8,\,1.25]\).
\item
  \textbf{Do not use \(t\)-tests.} Balance is a property of \emph{this
  sample}, not a population hypothesis --- and \(p\) confounds imbalance
  with sample size (Imai, King and Stuart 2008).
\end{itemize}
\end{frame}

\begin{frame}[fragile]{A balance plot: means are not enough}
\protect\phantomsection\label{a-balance-plot-means-are-not-enough}
\pandocbounded{\includegraphics[keepaspectratio]{03b_propensity_scores_files/figure-beamer/unnamed-chunk-1-1.pdf}}

\vspace{-0.3em}
\footnotesize

Selection here is linear in \texttt{Size}, \texttt{Leverage},
\texttt{ROA}, \texttt{Age} --- and \textbf{quadratic in
\texttt{Volatility}}. The score was fitted with linear terms only.
\end{frame}

\begin{frame}[fragile]{Reading the balance plot}
\protect\phantomsection\label{reading-the-balance-plot}
\begin{itemize}
\tightlist
\item
  \textbf{Left panel, the success.} Four covariates start badly
  imbalanced --- \texttt{Size} at \(0.65\), \texttt{Leverage} and
  \texttt{ROA} near \(0.49\). Compared \emph{within deciles of the
  score} all fall under \(0.03\). This is the plot referees want, and
  where most papers stop.
\item
  \textbf{Left panel, the trap.} \texttt{Volatility} looks
  \emph{perfect} from the start: \(\text{SMD} = -0.010\).
\item
  \textbf{Right panel.} Its variance ratio is \(\mathbf{2.27}\) ---
  treated firms have more than twice the spread. Extreme volatility
  drives selection in \emph{both} tails.
\item
  \textbf{Conditioning on the score makes it worse},
  \(2.27 \to \mathbf{2.80}\), because the fitted score never contained
  \(\texttt{vol}^2\). Add that one term and the ratio falls to
  \(\mathbf{1.07}\).
\item
  Judge the score by \textbf{balance achieved, not by fit} --- and
  balance means the whole distribution. Check squares and interactions,
  not just levels.
\end{itemize}
\end{frame}

\begin{frame}{Three things to do with the score}
\protect\phantomsection\label{three-things-to-do-with-the-score}
Once every unit has a \(\hat e(X_i)\), three standard uses:

\begin{enumerate}
\tightlist
\item
  \textbf{Diagnose.} Plot \(\hat e\) separately for treated and control
  units. This is the common-support check, and it costs nothing.
\item
  \textbf{Weight.} Reweight the sample to build a pseudo-experiment ---
  inverse probability weighting.
\item
  \textbf{Match.} Pair treated to control units with similar scores,
  then difference --- propensity score matching.
\end{enumerate}

\begin{itemize}
\tightlist
\item
  Busso, DiNardo and McCrary (2014): reweighting is competitive with the
  best matching estimators when overlap is good, whereas matching
  (e.g.~nearest neighbour with bias correction) can be more effective
  when overlap is poor.
\end{itemize}
\end{frame}

\begin{frame}{Visualizing overlap}
\protect\phantomsection\label{visualizing-overlap}
\pandocbounded{\includegraphics[keepaspectratio]{03b_propensity_scores_files/figure-beamer/unnamed-chunk-2-1.pdf}}

\vspace{-0.4em}

\footnotesize

Right panel: a treated mass at \(\hat e \approx 1\) with no controls,
and a control mass at \(\hat e \approx 0\) with no treated. Those units
have \textbf{no counterfactual}. Any estimate covering them is
extrapolation, not comparison.
\end{frame}

\begin{frame}{This plot belongs in every such paper}
\protect\phantomsection\label{this-plot-belongs-in-every-such-paper}
\begin{itemize}
\tightlist
\item
  Common support is verifiable, unlike unconfoundedness.
\item
  What to do when it looks like the right panel:

  \begin{itemize}
  \tightlist
  \item
    \textbf{Trim} extreme scores, or impose a caliper: drop units with
    \(\hat e \notin [0.1,\,0.9]\). A rule of thumb, but a defensible one
    (Crump, Hotz, Imbens and Mitnik, 2009).
  \item
    \textbf{Report the change in estimand.} After trimming you describe
    a subpopulation. Say how many observations went and how the treated
    sample shifted.
  \item
    Consider retreating to the \(ATT\), which only needs overlap in one
    direction.
  \end{itemize}
\item
  What \emph{not} to do: report the estimate and stay silent. Regression
  will extrapolate into a region with no data and give a standard error
  reflecting none of it.
\end{itemize}
\end{frame}

\begin{frame}{Where the weighting idea comes from}
\protect\phantomsection\label{where-the-weighting-idea-comes-from}
\begin{itemize}
\tightlist
\item
  The idea predates causal inference. \textbf{Horvitz and Thompson
  (1952)}, from survey sampling.
\item
  Among firms with \(\hat e = 0.1\), only one in ten was treated. The
  treated arm under-represents this type tenfold, so each such treated
  firm stands in for ten: weight it by \(1/\hat e = 10\). A treated firm
  with \(\hat e = 0.9\) is typical of its type: weight
  \(1/0.9 \approx 1.1\).
\item
  Do this in \textbf{both arms} --- treated by \(1/\hat e\), control by
  \(1/(1-\hat e)\) --- and you build a \textbf{pseudo-population} in
  which treatment no longer depends on \(X\). Selection is undone by
  \emph{recounting}, not by discarding or stratifying.
\item
  The weight depends on \textbf{which arm} a unit is in and on
  \textbf{which estimand} you want. And the danger is already visible: a
  score near \(0\) or \(1\) produces an enormous weight.
\end{itemize}
\end{frame}

\begin{frame}{Inverse probability weighting}
\protect\phantomsection\label{inverse-probability-weighting}
Create a balanced pseudo-population by inverse probability weighting.
Write \(\hat e_i\) for \(\hat e(X_i)\) and \(\pi = \Pr(D{=}1)\).

\small

\[
\begin{aligned}
ATE \;&=\; E\!\left[Y^1 - Y^0\right] = E\!\left[\frac{D\,Y}{e(X)}\right] - E\!\left[\frac{(1-D)\,Y}{1 - e(X)}\right] \\
ATT \;&=\; E\!\left[Y^1 - Y^0 \middle| D = 1 \right] =\frac{1}{\pi}\,\left(E\!\left[DY\right] - E\!\left[\frac{e(X)(1-D)Y}{1 - e(X)}\right]\right) \\[1em]
\widehat{ATE}_{IPW} \;&=\; \frac{1}{n}\sum_{i=1}^{n}\left[\frac{D_i}{\hat e_i} - \frac{1 - D_i}{1 - \hat e_i}\right] Y_i \\
\widehat{ATT}_{IPW} \;&=\; \frac{1}{n_T}\sum_{i=1}^{n}\left[D_i - (1 - D_i)\,\frac{\hat e_i}{1 - \hat e_i}\right] Y_i
\end{aligned}
\] \normalsize

\begin{itemize}
\tightlist
\item
  Both are \textbf{Horvitz--Thompson (1952)} forms.
\end{itemize}
\end{frame}

\begin{frame}{Why \emph{that} weight for the ATE?}
\protect\phantomsection\label{why-that-weight-for-the-ate}
We prove \(E\!\left[\frac{D\,Y}{e(X)}\right] = E[Y^1]\). The \(E[Y^0]\)
half is symmetric.

\scriptsize

\[
\begin{aligned}
E\!\left[ \frac{D \cdot Y}{e(X)} \right]
  &= E\!\left[ \frac{D \cdot Y^1}{e(X)} \right]
     && \text{switching equation (consistency): } D \cdot Y = D \cdot Y^1 \\
  &= E_X\!\left[ E\!\left( \frac{D \cdot Y^1}{e(X)} \,\Big|\, X \right) \right]
     && \text{law of total expectation} \\
  &= E_X\!\left[ \frac{1}{e(X)} \cdot E\big(D \cdot Y^1 \mid X\big) \right]
     && e(X) \text{ is fixed given } X \\
  &= E_X\!\left[ \frac{1}{e(X)} \cdot E(D \mid X) \cdot E\big(Y^1 \mid X\big) \right]
     && \text{unconfoundedness (exchangeability): } D \indep Y^1 \mid X \\
  &= E_X\!\left[ \frac{1}{e(X)} \cdot e(X) \cdot E\big(Y^1 \mid X\big) \right]
     && \text{by definition, } E(D \mid X) = e(X) \\
  &= E_X\big[ E(Y^1 \mid X) \big] \;=\; E[Y^1]
     && \text{cancel; law of total expectation}
\end{aligned}
\]
\end{frame}

\begin{frame}{Why \emph{that} weight for the ATT?}
\protect\phantomsection\label{why-that-weight-for-the-att}
\scriptsize

\(E[Y^1 \mid D{=}1] = E[Y \mid D{=}1] = \frac{E[DY]}{E[D]}\). \[
\begin{aligned}
E\big[Y^0 \mid D{=}1\big]
  &= E_X\big[\, E(Y^0 \mid X, D{=}1) \,\big|\, D{=}1 \,\big]
     && \text{law of total expectation, over } f(x \mid D{=}1) \\
  &= E_X\big[\, E(Y^0 \mid X, D{=}0) \,\big|\, D{=}1 \,\big]
     && \text{unconfoundedness --- swap the arm} \\
  &= E_X\!\left[\, E(Y^0 \mid X, D{=}0)\cdot\frac{e(X)}{\pi} \,\right]
     && \text{Bayes: } f(x \mid D{=}1) = e(x)f(x)/\pi \\
  &= E_X\!\left[\, \frac{E\big[(1-D)\,Y \mid X\big]}{1 - e(X)}\cdot\frac{e(X)}{\pi} \,\right]
     && E[(1{-}D)Y \mid X] = \big(1{-}e(X)\big)E[Y^0 \mid X, D{=}0] \\
  &= E\!\left[\, \frac{(1 - D)\,Y\,e(X)}{\big(1 - e(X)\big)\,\pi} \,\right]
     && \text{law of total expectation}
\end{aligned}
\] \normalsize

Line 3 converts an average over the \emph{treated} covariate
distribution into one over the \emph{whole} population.
\end{frame}

\begin{frame}{Ten banks, split by size}
\protect\phantomsection\label{ten-banks-split-by-size}
Lecture 1's capital injections, now with a covariate: bank
\textbf{size}.

\begin{center}
\scriptsize
\begin{tabular}{crcrrrr}
\toprule
Bank & Size & $D$ & $Y^0$ & $Y^1$ & $\tau$ & observed $Y$ \\
\midrule
$1$  & Small & $1$ & $10$ & $12$ & $2$ & $\mathbf{12}$ \\
$2$  & Small & $0$ & $8$  & $10$ & $2$ & $\mathbf{8}$  \\
$3$  & Small & $0$ & $9$  & $11$ & $2$ & $\mathbf{9}$  \\
$4$  & Small & $0$ & $11$ & $13$ & $2$ & $\mathbf{11}$ \\
$5$  & Small & $0$ & $12$ & $14$ & $2$ & $\mathbf{12}$ \\
\midrule
$6$  & Large & $1$ & $18$ & $24$ & $6$ & $\mathbf{24}$ \\
$7$  & Large & $1$ & $19$ & $25$ & $6$ & $\mathbf{25}$ \\
$8$  & Large & $1$ & $21$ & $27$ & $6$ & $\mathbf{27}$ \\
$9$  & Large & $1$ & $22$ & $28$ & $6$ & $\mathbf{28}$ \\
$10$ & Large & $0$ & $20$ & $26$ & $6$ & $\mathbf{20}$ \\
\bottomrule
\end{tabular}
\end{center}

\begin{itemize}
\tightlist
\item
  Small banks are \textbf{rarely} rescued, \(1\) of \(5\); large banks
  \textbf{usually} are, \(4\) of \(5\). So \(e(\text{Small}) = 0.2\) and
  \(e(\text{Large}) = 0.8\).
\item
  Small banks gain \(2\), large banks gain \(6\). Truth: \(ATE = 4\),
  \(ATT = 5.2\).
\item
  Only the \textbf{bold} column is observed. SDO = \(\mathbf{11.2}\).
\end{itemize}
\end{frame}

\begin{frame}{The ATE rebuilds the population}
\protect\phantomsection\label{the-ate-rebuilds-the-population}
\begin{center}
\small
\begin{tabular}{lrrrrrr}
\toprule
Cell & $N$ & $e$ & injected & $\times\,1/e$ & controls & $\times\,1/(1-e)$ \\
\midrule
Small & $5$ & $0.2$ & $1$ & $\times 5 = \mathbf{5}$    & $4$ & $\times 1.25 = \mathbf{5}$ \\
Large & $5$ & $0.8$ & $4$ & $\times 1.25 = \mathbf{5}$ & $1$ & $\times 5 = \mathbf{5}$ \\
\midrule
Total & $10$ & & $5$ & $\mathbf{10}$ & $5$ & $\mathbf{10}$ \\
\bottomrule
\end{tabular}
\end{center}

\begin{itemize}
\tightlist
\item
  Bank \(1\) is the \textbf{only} small bank rescued, so it must stand
  in for all five. Its weight is \(1/0.2 = 5\).
\item
  Every weighted entry comes back to \(5\): each arm is blown up to the
  \textbf{full population of ten}. That is the \emph{pseudo-population},
  and in it \(D \indep X\) by construction.
\item
  Before weighting, small banks are \(20\%\) of the injected sample and
  \(80\%\) of controls. After weighting, \(50\%\) of both --- the
  imbalance causing the bias is gone.
\end{itemize}
\end{frame}

\begin{frame}{The ATT reweights controls only}
\protect\phantomsection\label{the-att-reweights-controls-only}
The \(ATT\) wants the \textbf{treated} distribution. The four large and
one small rescued banks are already a sample from it, so they keep
weight \(1\). Only the controls move.

\begin{center}
\small
\begin{tabular}{lrrrr}
\toprule
Cell & controls & $\times\ e/(1-e)$ & $=$ & injected \\
\midrule
Small & $4$ & $0.25$ & $\mathbf{1}$ & $1$ \\
Large & $1$ & $4.00$ & $\mathbf{4}$ & $4$ \\
\bottomrule
\end{tabular}
\end{center}

\begin{itemize}
\tightlist
\item
  Multiplying the controls by \(e/(1-e)\) makes their count match the
  injected count \textbf{exactly} --- last column.
\item
  Bank \(10\) is the lone large control and must cover \textbf{four}
  rescued banks, so it carries weight \(4\).
\item
  If e were 1 --- all five large banks rescued, there is no large
  control to reweight and create counterfactuals: \textbf{this is where
  the \(ATT\) fails}, and why it needs only \(e(X) < 1\) while the
  \(ATE\) needs \(0 < e(X) < 1\).
\end{itemize}
\end{frame}

\begin{frame}{Does it work?}
\protect\phantomsection\label{does-it-work}
\[
\begin{aligned}
\widehat{E[Y^1]} &= \tfrac{1}{10}\!\left[\tfrac{12}{0.2} + \tfrac{24+25+27+28}{0.8}\right] = \tfrac{60 + 130}{10} = 19 \\
\widehat{E[Y^0]} &= \tfrac{1}{10}\!\left[\tfrac{8+9+11+12}{0.8} + \tfrac{20}{0.2}\right] = \tfrac{50 + 100}{10} = 15 \\[0.3em]
\widehat{ATE}\; &= \widehat{E[Y^1]} - \widehat{E[Y^0]} \;=\; 19 - 15 \;=\; \mathbf{4} \\[0.3em]
\widehat{ATT}\; &= \underbrace{\tfrac{12 + 24+25+27+28}{5}}_{\text{injected mean}\;=\;23.2} \;-\; \underbrace{\tfrac{(8+9+11+12)(0.25) + 20(4)}{4(0.25) + 1(4)}}_{\text{controls}\,\times\,\text{odds}\;=\;18.0} \;=\; \mathbf{5.2}
\end{aligned}
\]

\begin{itemize}
\tightlist
\item
  Truth: \(ATE = 4\), \(ATT = 5.2\). \textbf{Both exactly recovered.}
  The naive \(SDO\) is \(23.2 - 12 = \mathbf{11.2}\), nearly three times
  the \(ATE\).
\item
  \(ATT > ATE\) because large banks gain more, and large banks are where
  the injections went. \textbf{Selection on gains --- the two estimands
  correctly disagree.}
\end{itemize}
\end{frame}

\begin{frame}[fragile]{Two ways to divide}
\protect\phantomsection\label{two-ways-to-divide}
Both estimators use the same weights \(1/\hat e_i\). They differ
\textbf{only in the denominator}.

\vspace{-0.2em}

\[
\underbrace{\widehat{E[Y^1]}_{\text{HT}} = \frac{1}{N}\sum_i \frac{D_i Y_i}{\hat e_i}}_{\text{divide by } N}
\qquad\qquad
\underbrace{\widehat{E[Y^1]}_{\text{Hájek}} = \frac{1}{\sum_i D_i/\hat e_i}\sum_i \frac{D_i Y_i}{\hat e_i}}_{\text{divide by the weights you got}}
\]

\begin{itemize}
\tightlist
\item
  In the \textbf{population} they agree: \(E[D/e(X)] = 1\), so the
  weights sum to \(N\).
\item
  In a \textbf{sample} they do not. \(\sum_i D_i/\hat e_i\) is only an
  \emph{estimate} of \(N\), and it misses.
\item
  Hájek's weights sum to one by construction, so it is always a genuine
  \textbf{weighted average of outcomes you actually observed}. HT
  carries no such guarantee.
\item
  The default in Stata's \texttt{teffects\ ipw} and R's
  \texttt{WeightIt}.
\end{itemize}
\end{frame}

\begin{frame}{Example 1}
\protect\phantomsection\label{example-1}
Four rescued banks. \textbf{Every one has lending growth \(Y = 10\)}, so
the answer must be \(10\) for \(\widehat{E[Y^1]}\). Here \(N = 8\)
(including control with \(D=0\)).

\begin{center}
\footnotesize
\begin{tabular}{ccrc}
\toprule
Bank & $\hat e$ & weight $1/\hat e$ & $Y$ \\
\midrule
A & $0.50$ & $2$ & $10$ \\
B & $0.50$ & $2$ & $10$ \\
C & $0.50$ & $2$ & $10$ \\
D & $0.25$ & $4$ & $10$ \\
\midrule
sum & & $\mathbf{10}$ & \\
\bottomrule
\end{tabular}
\end{center}

\vspace{-0.3em}
\small

\[
\text{HT} = \frac{20+20+20+40}{8} = \frac{100}{8} = \mathbf{12.5},
\qquad
\text{Hájek} = \frac{100}{10} = \mathbf{10} .
\] \normalsize

\begin{itemize}
\tightlist
\item
  The weights sum to \(10\), not \(8\). HT divides by \(8\) anyway,
  returning \(10 \times \tfrac{10}{8}\).
\end{itemize}
\end{frame}

\begin{frame}{Example 2}
\protect\phantomsection\label{example-2}
\begin{center}
\footnotesize
\begin{tabular}{ccrrc}
\toprule
Bank & $\hat e$ & weight & share of weight & $Y$ \\
\midrule
A & $0.50$ & $2$  & $3.6\%$ & $10$ \\
B & $0.50$ & $2$  & $3.6\%$ & $10$ \\
C & $0.50$ & $2$  & $3.6\%$ & $10$ \\
D & $\mathbf{0.02}$ & $\mathbf{50}$ & $\mathbf{89.3\%}$ & $24$ \\
\midrule
sum & & $56$ & $100\%$ & \\
\bottomrule
\end{tabular}
\end{center}

\vspace{-0.3em}

\[
\text{HT} = \frac{1260}{8} = \mathbf{157.5},
\qquad
\text{Hájek} = \frac{1260}{56} = \mathbf{22.5} .
\]

\begin{itemize}
\tightlist
\item
  D carries \(\mathbf{89\%}\) of the weight.
\item
  \textbf{Normalization is necessary, not sufficient.} Check overlap
  separately: trim, or retreat to the \(ATT\) when the trouble is at
  \(\hat e \approx 0\).
\end{itemize}
\end{frame}

\begin{frame}{Propensity score matching}
\protect\phantomsection\label{propensity-score-matching}
\begin{itemize}
\tightlist
\item
  Match each treated unit to controls with a \textbf{similar score},
  then difference.

  \begin{itemize}
  \tightlist
  \item
    \textbf{Caliper} matching: refuse any match farther than some
    tolerance (say \(0.001\) on the score), \emph{discarding} the
    treated unit if none qualifies.
  \item
    \textbf{Radius} matching: average all controls inside the tolerance.
  \end{itemize}
\item
  Two problems with both:

  \begin{itemize}
  \tightlist
  \item
    The tolerance is \textbf{chosen by the researcher}
  \item
    Dropping treated units \textbf{changes the estimand}: the \(ATT\)
    summarised the \emph{original} treated sample.
  \end{itemize}
\item
  King and Nielsen (2019): two units can share a score,
  \(e(X_i) = e(X_j)\), with opposite covariate profiles
  \(X_i \neq X_j\). Matching on the scalar score can \emph{increase}
  local imbalance relative to multivariate distance matching.
\end{itemize}
\end{frame}

\begin{frame}{What the score does and does not buy}
\protect\phantomsection\label{what-the-score-does-and-does-not-buy}
\begin{itemize}
\tightlist
\item
  It buys \textbf{dimension reduction}, not identification. If \(X\)
  omits a confounder, \(e(X)\) omits it too.
\item
  Regression hides missing common support behind linear extrapolation;
  plotting \(\hat e(X)\) exposes it.
\item
  A high pseudo-\(R^2\) means covariates nearly predict treatment
  (\(e(X) \to 0\) or \(1\)): positivity breaks and IPW variance
  explodes. A low one means treatment is nearly independent of \(X\) ---
  close to an RCT on observables.
\item
  IPW shifts the modelling assumption from the outcome equation
  \(E[Y \mid D, X]\) to the assignment equation \(P(D{=}1 \mid X)\).
\end{itemize}
\end{frame}

\section{2. Regression}\label{regression}

\begin{frame}{Crossing to the other side of the map}
\protect\phantomsection\label{crossing-to-the-other-side-of-the-map}
\begin{itemize}
\tightlist
\item
  Everything so far has been an \textbf{interpolation} method:
  subclassification, matching, the propensity score. Common support is
  the binding constraint.
\item
  Regression sits on the other side: it needs no common support, so
  \textbf{the curse of dimensionality does not arise}.
\item
  It buys that with a \textbf{functional form}: OLS imputes the
  counterfactual by \textbf{extrapolating} into regions with no
  comparable units.
\item
  The trade: matching needs overlap; regression needs a correctly
  specified outcome model.
\item
  Both still need unconfoundedness. Nothing in this section relaxes it.
\end{itemize}
\end{frame}

\begin{frame}{Matching vs.~regression}
\protect\phantomsection\label{matching-vs.-regression}
\pandocbounded{\includegraphics[keepaspectratio]{03b_propensity_scores_files/figure-beamer/unnamed-chunk-5-1.pdf}}
\end{frame}

\begin{frame}{What exactly is that coefficient?}
\protect\phantomsection\label{what-exactly-is-that-coefficient}
\begin{itemize}
\tightlist
\item
  In practice you will estimate the effect by running \[
  Y_i = \alpha + \delta D_i + \gamma X_i + \varepsilon_i .
  \]
\item
  What \emph{is} \(\hat\delta\)? It is generally \textbf{neither the
  \(ATE\) nor the \(ATT\)}.
\item
  Two questions come first:

  \begin{enumerate}
  \tightlist
  \item
    \textbf{Which variation in \(D\)} does OLS use, once \(X\) is in the
    regression?
  \item
    \textbf{How does it weight} the units that supply that variation?
  \end{enumerate}
\item
  \textbf{Frisch--Waugh--Lovell} answers the first: OLS uses only the
  part of \(D\) that \(X\) cannot predict. The weighting result then
  follows from it.
\end{itemize}
\end{frame}

\begin{frame}{Coefficients as scaled covariances}
\protect\phantomsection\label{coefficients-as-scaled-covariances}
\begin{itemize}
\tightlist
\item
  With one regressor, \[
  Y_i = \beta_0 + \beta_1 X_i + u_i,
  \qquad
  \hat\beta_1 = \frac{\widehat{Cov}(Y_i, X_i)}{\widehat{Var}(X_i)} .
  \]
\item
  The coefficient is the slope of the best-fit line: a
  \textbf{covariance scaled by a variance}.
\item
  Now add covariates: \[
  Y_i = \beta_0 + \beta_1 X_{1i} + \beta_2 X_{2i} + \cdots + \beta_K X_{Ki} + \varepsilon_i .
  \]
\item
  \(\hat\beta_1\) is no longer
  \(\widehat{Cov}(Y, X_1)/\widehat{Var}(X_1)\). \textbf{Frisch and Waugh
  (1933); Lovell (1963)}: it is still a scaled covariance --- of
  \emph{residualized} variables.
\end{itemize}
\end{frame}

\begin{frame}{Frisch--Waugh--Lovell in three steps}
\protect\phantomsection\label{frischwaughlovell-in-three-steps}
\textbf{Step 1. Auxiliary regression}: \[
X_{1i} = \gamma_0 + \gamma_2 X_{2i} + \cdots + \gamma_K X_{Ki} + f_i .
\] \(\hat X_{1i}\) is the part of \(X_{1i}\) the other covariates
explain.

\textbf{Step 2. Residualization}: \[
\tilde X_{1i} = X_{1i} - \hat X_{1i} .
\] The variation in \(X_{1i}\) that is \emph{uncorrelated} with
everything else --- \textbf{partialling out}.

\textbf{Step 3. Regress the outcome on the residual.} \[
Y_i = \alpha + \beta_1 \tilde X_{1i} + \varepsilon_i,
\qquad
\hat\beta_1 = \frac{\widehat{Cov}(Y_i, \tilde X_{1i})}{\widehat{Var}(\tilde X_{1i})} .
\]

\begin{itemize}
\tightlist
\item
  \(\hat\beta_1\) is \textbf{identical} to the coefficient on \(X_{1i}\)
  in the full regression.
\end{itemize}
\end{frame}

\begin{frame}{FWL is an identity, not an approximation}
\protect\phantomsection\label{fwl-is-an-identity-not-an-approximation}
Simulation: \(4{,}000\) workers, earnings on a treatment indicator
\(d\), age, and high-school GPA, with a true effect of \(\$1{,}500\).

\begin{center}
\footnotesize
\begin{tabular}{lcc}
\toprule
 & (1) Multivariate OLS & (2) FWL \\
\midrule
$d$   & $1503.201$ & $1503.201$ \\
      & $(8.249)$  & $(8.287)$ \\
age   & $1.173$    & --- \\
      & $(2.023)$  & \\
gpa   & $9.081$    & --- \\
      & $(4.615)$  & \\
$N$   & $4000$     & $4000$ \\
\bottomrule
\end{tabular}
\end{center}

\begin{itemize}
\tightlist
\item
  The estimate sits \(\$3.20\) above the true \(\$1{,}500\) --- \(0.39\)
  standard errors, i.e.~sampling noise.
\item
  The standard errors do \emph{not} agree, and that is not sampling
  noise.
\end{itemize}
\end{frame}

\begin{frame}{Residualize both sides}
\protect\phantomsection\label{residualize-both-sides}
\begin{itemize}
\tightlist
\item
  Another version residualizes the outcome too: \[
  \begin{aligned}
  \tilde Y_i \;=\; Y_i \;-\; \hat Y_i &\quad\text{from}\quad Y_i = \lambda_0 + \lambda_2 X_{2i} + \cdots + \lambda_K X_{Ki} + v_i,\\
  &\quad\text{then}\quad \tilde Y_i = \beta_1 \tilde X_{1i} + \varepsilon_i.
  \end{aligned}
  \]
\item
  Both versions return \textbf{exactly the same \(\hat\beta_1\)}.
\item
  The double-residual version buys one more thing: \textbf{its residuals
  are those of the full regression}.
\end{itemize}
\end{frame}

\begin{frame}{Three regressions, one coefficient}
\protect\phantomsection\label{three-regressions-one-coefficient}
\begin{center}
\footnotesize
\begin{tabular}{lccc}
\toprule
 & (1) full OLS & (2) single, $Y \sim \tilde D$ & (3) double, $\tilde Y \sim \tilde D$ \\
\midrule
\multicolumn{4}{l}{\emph{original:} the controls explain $0.1\%$ of earnings} \\
$d$ & $1503.201$ & $1503.201$ & $1503.201$ \\
    & $(8.249)$  & $(8.287)$  & $(8.247)$ \\
\midrule
\multicolumn{4}{l}{\emph{expanded:} the controls explain $93\%$} \\
$d$ & $1503.201$ & $1503.201$ & $1503.201$ \\
    & $(8.249)$  & $(\mathbf{91.493})$ & $(8.247)$ \\
\bottomrule
\end{tabular}
\end{center}

\begin{itemize}
\tightlist
\item
  \begin{enumerate}
  [(1)]
  \tightlist
  \item
    and (3) differ only in degrees of freedom (\(n-k-1\) vs.~\(n-2\));
    the residuals are identical.
  \end{enumerate}
\item
  \begin{enumerate}
  [(1)]
  \setcounter{enumi}{1}
  \tightlist
  \item
    residualizes only \(D\). The part of \(Y\) that age and GPA explain
    goes to the error term. Harmless at \(R^2 = 0.1\%\); fatal at
    \(93\%\).
  \end{enumerate}
\end{itemize}
\end{frame}

\begin{frame}{Why FWL is worth knowing}
\protect\phantomsection\label{why-fwl-is-worth-knowing}
\begin{itemize}
\tightlist
\item
  \textbf{It tells you what variation identifies your coefficient.}
  \(\hat\beta_1\) uses only the part of \(X_1\) that the controls cannot
  explain.
\item
  \textbf{It makes the estimate visualizable.} Residualize \emph{both}
  sides and scatter \(\tilde Y\) against \(\tilde X_1\): a multivariate
  coefficient becomes a picture whose slope is \(\hat\beta_1\).
\item
  \textbf{It is the engine behind much else.} Fixed effects, double
  machine learning and the weighting results below are all FWL in
  disguise.
\item
  Angrist and Pischke (2009) call this \textbf{regression anatomy}, and
  it is the single most useful algebraic fact about OLS.
\end{itemize}
\end{frame}

\begin{frame}{Exogeneity, and the canonical OLS model}
\protect\phantomsection\label{exogeneity-and-the-canonical-ols-model}
\begin{itemize}
\tightlist
\item
  The \textbf{canonical OLS model}: \[
  Y_i = \alpha + \delta D_i + \gamma X_i + \varepsilon_i.
  \]
\item
  \(E[\varepsilon_i \mid D_i, X_i] = 0\) means ``\(D\) is exogenous.''
  It bundles three assumptions:

  \begin{itemize}
  \tightlist
  \item
    \textbf{Unconfoundedness}: \(E[Y^0 \mid D, X] = E[Y^0 \mid X]\).
  \item
    \textbf{Linearity}: \(E[Y^0 \mid X] = \alpha + \gamma X\).
  \item
    \textbf{Constant effect}: \(\tau_i = \delta\) for every \(i\).
  \end{itemize}
\item
  We take \(X\) to be pre-treatment throughout. If \(D\) moves \(X\),
  \(X\) is an outcome, not a control (Lecture 2).
\end{itemize}
\end{frame}

\begin{frame}{When the three assumptions fail}
\protect\phantomsection\label{when-the-three-assumptions-fail}
Each failure pushes something into \(\varepsilon_i\) that is correlated
with \(D_i\).

\begin{itemize}
\tightlist
\item
  \textbf{Linearity fails:}
  \(E[Y^0 \mid X] = g(X) \neq \alpha + \gamma X\). The unmodeled part
  \(g(X) - \alpha - \gamma X\) enters \(\varepsilon_i\) and is
  correlated with \(D\) through \(X\).

  \begin{itemize}
  \tightlist
  \item
    \emph{Fix:} flexible controls --- polynomials, splines, category
    dummies
  \end{itemize}
\item
  \textbf{Constant effect fails:} \(\tau_i \neq \delta\). The term
  \(D_i(\tau_i - \bar\tau)\) enters \(\varepsilon_i\).
  \(\hat\delta_{OLS}\) becomes a variance-weighted average of
  \(\tau_i\), overweighting cells with high \(Var(D \mid X)\), and can
  differ sharply from the ATE.

  \begin{itemize}
  \tightlist
  \item
    \emph{Fix:} interact \(D_i\) with covariates
  \end{itemize}
\item
  \textbf{Unconfoundedness fails:}
  \(E[Y^0 \mid D, X] \neq E[Y^0 \mid X]\), i.e.~selection on
  unobservables: \(E[Y^0 \mid D{=}1, X] \neq E[Y^0 \mid D{=}0, X]\).

  \begin{itemize}
  \tightlist
  \item
    \emph{Fix:} change the design!
  \end{itemize}
\end{itemize}
\end{frame}

\begin{frame}{OLS as a variance-weighted average}
\protect\phantomsection\label{ols-as-a-variance-weighted-average}
\begin{center}
\fbox{\;\textbf{Claim}\quad With $X$ saturated and unconfoundedness, $\;\hat\delta_{OLS} \to E\big[\sigma^2_D(X)\,\tau(X)\big] \,\big/\, E\big[\sigma^2_D(X)\big]$\;}
\end{center}

\scriptsize

\emph{Saturated}: a dummy per value of \(X\), so \(D\) on the controls
gives \(e(X)\) exactly. Let \(\tau(X) = E[\tau_i \mid X]\),
\(\sigma^2_D(X) = Var(D \mid X) = e(1-e)\).

\textbf{Step 1 (FWL).} \(\tilde D = D - e(X)\) has mean zero, so
\(\delta_{OLS} = Cov(Y, \tilde D)/Var(\tilde D) = E\big[Y(D - e)\big] / E\big[(D - e)^2\big]\).

\textbf{Step 2 (within a cell).} The denominator is
\(E\big[(D - e)^2 \mid X\big] = \sigma^2_D(X)\). For the numerator,
substitute
\(E[Y \mid X] = e\,E[Y \mid D{=}1,X] + (1-e)\,E[Y \mid D{=}0,X]\): \[
E\big[Y(D - e) \mid X\big] = e\,E[Y \mid D{=}1,X] - e\,E[Y \mid X]
= e(1-e)\big(E[Y \mid D{=}1,X] - E[Y \mid D{=}0,X]\big) = \sigma^2_D(X)\,\tau(X),
\] the last step by \textbf{unconfoundedness}. \textbf{Step 3.} Iterate
expectations over \(X\) in numerator and denominator. \(\blacksquare\)

Weights depend on \(i\) only through \(X_i\): largest where
\(e(X) = \tfrac12\), \textbf{zero} where \(e(X) \in \{0,1\}\). No
functional form for \(Y\) was used --- saturation absorbed it.

\textbf{Słoczyński (2022).} If \(X\) enters \emph{linearly} instead,
with \(Y^0, Y^1\) linear in \(X\):
\(\;\delta_{OLS} = w_1\,ATT + (1 - w_1)\,ATU\), where \[
w_1 = \frac{(1-\rho)\,V[p(X) \mid D{=}0]}{\rho\,V[p(X) \mid D{=}1] + (1-\rho)\,V[p(X) \mid D{=}0]}, \qquad \rho = \Pr(D{=}1), \quad p(X) = \text{linear projection of } D \text{ on } X.
\] Equal variances give \(w_1 = 1 - \rho\): \textbf{the \(ATT\) gets the
control share} --- the smaller group gets the larger weight.
\end{frame}

\begin{frame}[fragile]{Regression adjustment}
\protect\phantomsection\label{regression-adjustment}
\begin{itemize}
\item
  Interact treatment with every covariate: \[
  Y_i = \alpha + \delta D_i + \gamma_0 X_i + \lambda\, D_i X_i + \varepsilon_i,
  \qquad
  \widehat{ATT} = \hat\delta + \hat\lambda \bar X_1, \quad
  \widehat{ATE} = \hat\delta + \hat\lambda \bar X .
  \]
\item
  Or center first: interact with \((X_i - \bar X_1)\) and read the
  \(ATT\) off \(\hat\delta\); with \((X_i - \bar X)\) for the \(ATE\).
  Same coefficients, different covariate distribution.
\item
  Equivalently: fit \(Y\) on \(X\) separately in each arm. For the
  \(ATT\), impute \(Y^0\) for treated units from the control fit and
  average the gap over the treated. For the \(ATE\), also impute \(Y^1\)
  for controls from the treated fit and average over everyone.
  \texttt{teffects\ ra} does both and gets the SEs right.
\end{itemize}
\end{frame}

\section{3. The NSW Demonstration}\label{the-nsw-demonstration}

\begin{frame}{The 1970s empirical crisis}
\protect\phantomsection\label{the-1970s-empirical-crisis}
\begin{itemize}
\tightlist
\item
  The methods in this lecture were not adopted because they were
  elegant. They were adopted after a crisis of confidence in
  observational work.
\item
  Job-training programs in the 1970s were evaluated observationally, and
  the estimates were all over the place --- depending, apparently, on
  the analyst.
\item
  The \textbf{National Supported Work (NSW) Demonstration} was
  different: it \emph{randomized} participation. So the experimental
  effect was known.
\item
  That created a \textbf{benchmark}: could observational methods,
  applied to the same treated units, recover the experimental answer?
\end{itemize}
\end{frame}

\begin{frame}{What LaLonde did}
\protect\phantomsection\label{what-lalonde-did}
\begin{itemize}
\tightlist
\item
  LaLonde (1986) ran the decisive test: take the NSW experiment,
  \textbf{throw away the experimental control group}, and replace it
  with comparison groups built from survey data (CPS and PSID).
\item
  He applied the standard observational toolkit and compared the results
  to the experimental benchmark of roughly \textbf{\$1,800} in annual
  earnings gains.
\item
  The estimates ranged from a fraction of the benchmark to \textbf{the
  wrong sign}, depending on how the comparison group was built.
\item
  The paper was devastating, and a major reason the profession moved
  toward IV, RDD and DiD.
\end{itemize}
\end{frame}

\begin{frame}{Dehejia and Wahba: not so fast}
\protect\phantomsection\label{dehejia-and-wahba-not-so-fast}
\begin{itemize}
\tightlist
\item
  \textbf{Dehejia and Wahba (1999, 2002)} revisited LaLonde's data with
  propensity score methods.
\item
  Their finding: the observational estimates \emph{did} recover the
  experimental benchmark reasonably well --- \textbf{once common support
  was imposed}.
\item
  Much of LaLonde's failure was an \textbf{overlap} problem, not only
  confounding: the CPS and PSID groups held huge numbers of people
  utterly unlike NSW participants.
\item
  Discard the non-comparable units and the methods perform far better.
\end{itemize}
\end{frame}

\begin{frame}{Smith and Todd: not so fast either}
\protect\phantomsection\label{smith-and-todd-not-so-fast-either}
\begin{itemize}
\tightlist
\item
  \textbf{Smith and Todd (2005)} re-examined the Dehejia--Wahba result
  and found it fragile.
\item
  The success was \textbf{sensitive to the subsample} analysed and to
  the \textbf{propensity score specification} chosen. Reasonable
  alternative choices undid it.
\item
  Their charge is uncomfortable and worth stating plainly: the
  specification that worked may have been found \emph{because} the
  experimental answer was already known.
\item
  You cannot run that search in real research: there is no benchmark to
  check against --- that is the point of doing the study.
\end{itemize}
\end{frame}

\begin{frame}{The honest summary}
\protect\phantomsection\label{the-honest-summary}
\begin{itemize}
\tightlist
\item
  Selection on observables \textbf{can} work --- when overlap is genuine
  and the confounders are truly measured.
\item
  Neither condition is verifiable in the setting where you need them:

  \begin{itemize}
  \tightlist
  \item
    Unconfoundedness is \textbf{untestable}, always.
  \item
    Common support is \textbf{checkable}.
  \end{itemize}
\item
  Do the check you can do, and argue explicitly for the assumption you
  cannot.
\item
  \textbf{A useful exercise}: the NSW data are freely available. Apply
  every method here and compare to the \$1,800 benchmark --- the fastest
  way to build instincts about when these methods hold up.
\end{itemize}
\end{frame}

\section{Wrapping Up}\label{wrapping-up}

\begin{frame}{Summary {[}Part 1{]}: the score and what it does}
\protect\phantomsection\label{summary-part-1-the-score-and-what-it-does}
\begin{itemize}
\tightlist
\item
  \textbf{The propensity score} reduces \(X\) to one number. Balance
  given \(e(X)\) is free; the theorem --- conditioning on \(e(X)\)
  suffices --- needs unconfoundedness. It buys dimension reduction,
  never \textbf{identification}.
\item
  \textbf{Judge the score by balance, not fit} --- and balance means the
  whole distribution.
\item
  \textbf{Plot the overlap.} Common support is the one assumption you
  can check. Trim if you must, and say how the estimand changed.
\item
  \textbf{IPW} reweights into a pseudo-population: \(1/\hat e\) and
  \(1/(1-\hat e)\) for the \(ATE\); \(\hat e/(1-\hat e)\) on controls
  for the \(ATT\). Normalize (Hájek), and watch the tails.
\item
  \textbf{Propensity score matching} is the weakest use: two units can
  share a score with opposite covariates, so pruning can \emph{increase}
  imbalance (King and Nielsen 2019).
\end{itemize}
\end{frame}

\begin{frame}{Summary {[}Part 2{]}: regression, and the verdict}
\protect\phantomsection\label{summary-part-2-regression-and-the-verdict}
\begin{itemize}
\tightlist
\item
  \textbf{The trade.} Matching needs overlap; regression needs a
  correctly specified outcome model. Both still need unconfoundedness.
\item
  \textbf{FWL.} OLS identifies \(\delta\) only from the part of \(D\)
  that \(X\) cannot predict.
\item
  \textbf{The canonical model} bundles unconfoundedness, linearity and
  constant effects. With heterogeneous effects it returns a
  variance-weighted average, heaviest where treatment is a coin flip ---
  generally neither the \(ATE\) nor the \(ATT\) (Angrist 1998;
  Słoczyński 2022).
\item
  \textbf{Regression adjustment} --- interacting \(D\) with \(X\) ---
  recovers both, but still leans on functional form.
\item
  \textbf{NSW.} These methods work when overlap is real and the
  confounders are measured --- and you can only ever verify the first.
\item
  Next lecture: designs that \emph{create} variation instead of assuming
  it away.
\end{itemize}
\end{frame}

\begin{frame}{References}
\protect\phantomsection\label{references}
\tiny

\begin{itemize}
\tightlist
\item
  Angrist, J. (1998). Estimating the Labor Market Impact of Voluntary
  Military Service Using Social Security Data on Military Applicants.
  \emph{Econometrica} 66, 249--288.
\item
  Angrist, J. and J.-S. Pischke (2009). \emph{Mostly Harmless
  Econometrics: An Empiricist's Companion}. Princeton University Press.
\item
  Busso, M., J. DiNardo and J. McCrary (2014). New Evidence on the
  Finite Sample Properties of Propensity Score Reweighting and Matching
  Estimators. \emph{Review of Economics and Statistics} 96, 885--897.
\item
  Crump, R., V. J. Hotz, G. Imbens and O. Mitnik (2009). Dealing with
  Limited Overlap in Estimation of Average Treatment Effects.
  \emph{Biometrika} 96, 187--199.
\item
  Cunningham, S. \emph{Causal Inference: The Remix}, Ch. 5 (§5.9 for
  regression).
\item
  Dehejia, R. and S. Wahba (1999). Causal Effects in Nonexperimental
  Studies. \emph{JASA} 94, 1053--1062.
\item
  Dehejia, R. and S. Wahba (2002). Propensity Score-Matching Methods for
  Nonexperimental Causal Studies. \emph{Review of Economics and
  Statistics} 84, 151--161.
\item
  Frisch, R. and F. Waugh (1933). Partial Time Regressions as Compared
  with Individual Trends. \emph{Econometrica} 1, 387--401.
\item
  Hájek, J. (1971). Comment on Basu. In \emph{Foundations of Statistical
  Inference}, Holt, Rinehart and Winston, 236.
\item
  Horvitz, D. and D. Thompson (1952). A Generalization of Sampling
  Without Replacement from a Finite Universe. \emph{JASA} 47, 663--685.
\item
  Imai, K., G. King and E. Stuart (2008). Misunderstandings between
  Experimentalists and Observationalists about Causal Inference.
  \emph{JRSS Series A} 171, 481--502.
\item
  King, G. and R. Nielsen (2019). Why Propensity Scores Should Not Be
  Used for Matching. \emph{Political Analysis} 27, 435--454.
\item
  LaLonde, R. (1986). Evaluating the Econometric Evaluations of Training
  Programs with Experimental Data. \emph{AER} 76, 604--620.
\item
  Lovell, M. (1963). Seasonal Adjustment of Economic Time Series and
  Multiple Regression Analysis. \emph{JASA} 58, 993--1010.
\item
  Rosenbaum, P. and D. Rubin (1983). The Central Role of the Propensity
  Score in Observational Studies for Causal Effects. \emph{Biometrika}
  70, 41--55.
\item
  Słoczyński, T. (2022). Interpreting OLS Estimands When Treatment
  Effects Are Heterogeneous: Smaller Groups Get Larger Weights.
  \emph{Review of Economics and Statistics} 104(3), 501--509.
\item
  Smith, J. and P. Todd (2005). Does Matching Overcome LaLonde's
  Critique of Nonexperimental Estimators? \emph{Journal of Econometrics}
  125, 305--353.
\end{itemize}
\end{frame}




\end{document}
